%!TEX root = ../../../main.tex

\chapter{Grundlagen}
\label{chap:grundlagen}

\paragraph{RAG.} Ein Sprachmodell kennt meine Rechnungen nicht. Bei Retrieval-Augmented Generation wird deshalb vor der Antwort gesucht, und die gefundenen Textstellen kommen mit in den Prompt \cite{lewis2020}. Gesucht wird meist über Embeddings, also Vektoren für Textblöcke und Frage. Gao et~al. unterscheiden einfaches RAG (\enquote{suchen, dann antworten}) von \emph{modularem} RAG aus austauschbaren Bausteinen wie Vorverarbeitung, Suche, Umformulieren, Neusortieren und Prüfen \cite{gao2023}. Diesen Gedanken nutze ich: Jeden Baustein messe ich einzeln und entferne ihn wieder, wenn er nichts bringt.

\paragraph{Agenten und Werkzeuge.} Ein Agent legt den Ablauf selbst fest. Beim ReAct-Muster wechselt das Modell zwischen Gedanke, Werkzeugaufruf und Beobachtung, bis es antworten kann \cite{yao2023}. Toolformer zeigt, dass Modelle lernen können, wann ein Werkzeug wie ein Rechner hilft \cite{schick2023}; gerade kleine Modelle verrechnen sich oft. Bei Self-RAG entscheidet das Modell selbst, ob es suchen muss, und prüft danach, ob seine Antwort von den Stellen gedeckt ist \cite{asai2024}. GraphRAG baut aus allen Texten einen Wissensgraphen und fasst Gruppen davon zusammen \cite{edge2024}. Das hilft bei Fragen über die ganze Sammlung, kostet beim Einlesen aber sehr viele Modellaufrufe (Abschnitt~\ref{sec:graphrag}).
